\begin{abstract}
We encode consecutive primes as residue classes modulo $30$, producing a sequence over 8 states $\{1,7,11,13,17,19,23,29\}$. We construct a first-order transition operator $P$ and compare the empirical two-step operator $P^{(2)}$ to the baseline $P^2$ via $\Delta = P^{(2)} - P^2$.

Using the first $N$ primes, $\Delta$ has low-rank structure. Low-rank approximations improve prediction of $P^{(2)}$ relative to $P^2$.

Synthetic controls produce distinct residual operators. Residue-class transitions of consecutive primes modulo $30$ are represented by the residual operator $\Delta = P^{(2)} - P^2$, which exhibits higher-order structure.
\end{abstract}